\section{Definition and basic properties of one-dimensional
formal group laws}

Let $R$ be a commutative ring throughtout this chapter.

\begin{definition}[{\cite[Definition 1.1]{hazewinkel1978formal}}]
    A \emph{one-dimensional formal group law over $R$} is a formal power
    series $F(X, Y) \in R[\![X,Y]\!]$ with the following properties:
    \begin{enumerate}
        \item $F(X,Y)$ is of the form
        \begin{equation}
            F(X,Y) = X + Y + \sum_{i, j\ge 1}c_{ij} X^i Y^j.
        \end{equation}
        \item $F(X,Y)$ satisfies the associativity condition, namely
        \begin{equation}
            F(F(X,Y),Z) = F(X,F(Y,Z)).
        \end{equation}
    \end{enumerate}
    For convention, we use formal group law to refer to one-dimensional formal group law.
    Furthermore, if $F(X,Y)$ satisfies $F(X,Y) = F(Y,X)$, then it is called
    a \emph{commutative formal group law}.
    \label{Defn formal group}
\end{definition}

\begin{remark}
    Under this definition, one can prove that $F(X,0) = X$ and $F(0,Y) = Y$, which are included
    as part of the definition in other references.

    To see this, let  $f(X) = F(X,0)$. Then we have
        \[F(F(X,0),0)= F(X,F(0,0))= F(X,0).\]
    which
    means that $f \circ f = f$. Composing both sides with $f^{-1}$ gives $f = \id$.
    This proves that $F(X,0) = X$. The equality $F(0,Y)= Y$ follows from an analogous argument.

    Therefore, given a formal group law $F(X,Y) \in R[\![X,Y]\!]$, it can be written as
    $F(X,Y)= X + YG_1(X,Y)$ and $F(X,Y) = Y + XG_2(X,Y)$,
    for some $G_1(X,Y), G_2(X,Y) \in R [\![X,Y]\!]$.
\end{remark}


\begin{example}
    Two examples of commutative formal group laws are:
    \begin{align}
        & \GG_a (X,Y) = X + Y \quad \text{(the additive formal group law)} \\
        & \GG_m (X,Y) = X + Y + XY \quad \text{(the multiplicative formal group law)}
    \end{align}
    %% TODO: explain why there are called additive or multiplicative formal group law.
\end{example}

\begin{lemma}
    Let $F(X, Y)$ be a formal group law over a ring $R$. Then there exists a power series
    $i(X) \in R [\![X]\!]$ such that $F(X, i(X)) = 0$ and $F(i(X),X) = 0$.
    In particular, the left inverse and right inverse are the same and the inverse is unique.
    \label{lem: additive inverse}
\end{lemma}

\begin{proof}
    The idea is to write $i(X) = \sum_{j= 1}^\infty b_j X ^ j$, and choose $b_j$ inductively so that
    $F(X,i(X)) = 0$. \par
    First, we choose $b_1 = -1$, so that $F(X, i(X)) = 0 \pmod{\deg 2}$. Assume that
    $b_1, \ldots, b_n$ are already chosen so that $F(X, i(X)) = 0 \pmod{\deg (n + 1)}$.
    Then choose $b_{n+1}$ to be minus the $(n+1)$-th coefficient of $F(X, i_n(X))$, where
    $i_n(X)= \sum_{j = 1}^n b_j X ^ j$. Therefore, $F(X, i_{n+1}(X)) = 0 \pmod{\deg (n + 2)}$,
    where $i_{n+1}(X) = i_n(X) + b_{n+1}X ^{n+1}$. Repeating this process, we obtain the desired
    series $i(X)$ with $F(X,i(X)) = 0$. We can use a similar process to define $i'(X)$ such that
    $F(i'(X), X) = 0$. \par
    Since $F(X,Y)$ is a formal group law, we have
    \[ i'(X) = F(i'(X), 0) = F(i'(X), F(X, i(X))) = F(F(i'(X), X), i(X)) =
    F(0,i(X)) = i(X)
    \]
    Therefore, the left inverse and right inverse of $X$ are the same.

    For the uniqueness, let $i_1(X), i_2(X) \in R[\![X]\!]$ such that
    $F(X,i_1(X)) = 0 = F(X, i_2(X))$, then
    \begin{align*}
        i_1(X) = F(i_1(X), 0) &= F(i_1(X), F(X, i_2(X))) \\
        &= F(F(i_1(X), X), i_2(X)) = F(0, i_2(X)) = i_2(X).
    \end{align*}

\end{proof}

\begin{remark}
    Let $S$ be commutative rings and $S$ is an $R$-algebra.
    A formal group law $F(X,Y)$ over $R$ defines a additive group structure on the set of
    multivariate power series in $S[\![X_I]\!]$ with nilpotent constant coefficient,
    where $I$ is an idex set.
    Moreover, let $S$ be a complete topological ring, $F(X,Y)$ defines a additive group
    structure on the set of topological nilpotents of $S$. Completeness of $S$ ensure that the
    evaluation of power series $F(X,Y)$ converges to unique element in $S$.

\end{remark}

\begin{definition}[{\cite[Definition 1.4]{hazewinkel1978formal}}]
    Let $F(X,Y)$ and $G(X,Y)$ be two formal group laws over a ring $R$.
    A \emph{homomorphism of formal group laws}
    $F(X,Y) \to G(X,Y)$ is a power series $\alpha (X) = \sum_{i=1}^{\infty}a_i X ^ i \in R[\![X]\!]$ such that
    \begin{equation}
        \alpha (F(X,Y))= G(\alpha(X), \alpha(Y))
    \end{equation}
    An \emph{endomorphism of formal group law} $F(X,Y)$ is a homomorphism $F(X,Y) \to F(X,Y)$.
    A homomorphism $\alpha(X)$ is an \emph{isomorphism} if there exists a homomorphism
    $\beta(X) : G(X,Y) \to F(X,Y)$ such that $\alpha (\beta(X)) = \beta (\alpha (X))= X$.
    \label{defn: hom of formal group law}
\end{definition}

\begin{lemma}
    Let $f(X) = \sum_{i=1}^{\infty} a_i X ^ i \in R[\![X]\!]$, assume that $a_1$ is invertible in
    $R$, then there is an
    unique power series $g(X) \in R[\![X]\!]$ such that $f(g(X)) = X = g(f(X))$.
    \label{lem: substInv}
\end{lemma}

\begin{proof}
    The idea is to write $g(X) = \sum_{i=1}^{\infty} b_i X ^ i$ and choose inductively $b_i$
    so that $f(g(X))= X$.

    First, we choose that $b_1 = a_1 ^{-1}$ so that $f(g(X)) = X \pmod{\deg 2}$. Assume that
    $b_1, \ldots, b_n$ is already chosen, we define $c_{n+1}$ be the $n+1$-th coefficient of
    $f(\sum_{i=1}^{n}b_i X ^ i)$ and $b_{n+1} = - a_1^{-1} c_{n+1}$. Then the $n+1$-th coefficient
    of $f(\sum_{i=1}^{n+1} b_i X ^ i)$ is $c_{n+1} + a_1 b_{n+1} = 0$. Repeating this process,
    we get a power series $g(X) \in R[\![X]\!]$ so that $f(g(X)) = X$.

    Secondly, since $b_1$ is invertible in $R$, then there is a power series $h(X) \in R[\![X]\!]$
    so that $g(h(X)) = X$. Precomposing $f$ at the both side of the equality $g(h(X)) = X$, then
    we have $f(g(h(X))) = f(X)$. Therefore, we conclude that $h(X) = f(X)$, i.e.
    $g(f(X)) = X$.

    Let $g_1(X)$ and $g_2(X)$ be two power series so that
    $f(g_i(X)) = X = g_i(f(X))$. Then we have that
    \[g_1(X) = g_1(f(g_2(X))) = g_2(X). \]
\end{proof}

\begin{lemma}
    The homomorphism $\alpha(X) : F(X,Y) \to G(X,Y) $ is an isomorphism if and only
    if $a_1$ is a unit in $R$, where $\alpha (X) = \sum_{i=1}^{\infty} a_i X ^ i $.
    \label{lem: hom_iff_coeff_unit}
\end{lemma}

\begin{proof}
    Necessity: given that $\alpha (X)$ is a power series with $a_1$ a unit, there exists a power series
    $\beta (X)$ such that $\alpha (\beta (X)) =\beta (\alpha (X)) = X$. We need to prove that
    this $\beta (X)$ is a formal group homomorphism from $G(X,Y)$ to $F(X,Y)$, namely
    $\beta (G(X,Y)) = F (\beta (X), \beta (Y))$. \par
    Since $\alpha (X)$ is a homomorphism from $F(X,Y)$ to $G(X,Y)$,
    \[ \alpha (F(\beta(X), \beta (Y))) = G(\alpha (\beta (X)), \alpha (\beta (Y))) = G(X,Y).\]
    The desired equality follows by composing both sides with $\beta (X)$.

    Sufficiency: assume that $\alpha(X)$ is an formal group isomorphism, then there
    is formal group homomorphism $\beta(X) : G(X) \to F(X)$ such that $\alpha(\beta(X)) = X$.
    If we denote that $\alpha(X) = \sum_{i=1}^{\infty} a_i X ^ i$ and
    $\beta(X) = \sum_{i=1}^{\infty}b_i X ^ i$, then the coefficient of $X$ on the
    $\alpha(\beta(X))$ is $a_1 b _1$. Therefore, we conclude that $a_1$ is invertible
    in $R$, since $a_1 b_1 = 1$.
\end{proof}

One of the most important endomorphisms of a commutative formal group is the $n$-series.

\begin{definition}[$n$-series, {\cite[Definition 1.4.5]{hazewinkel1978formal}}]
    Let $F(X,Y)$ be a commutative formal group law over $R$,
    for natural number $n$, we define the $n$-series $[n](X)\in R[\![X]\!]$ inductively
    as follows:
    \begin{enumerate}
        \item $[0](X) = 0$.
        \item $[n+ 1](X) = F([n](X), X)$.
    \end{enumerate}
    Moreover, $[-n](X) \coloneqq i ([n](X))$, where $i(X)$ is the power series in Lemma
    \ref{lem: additive inverse}.
\end{definition}


\begin{proposition}
    For any natural number $n$, the $n$-series $[n](X)$ is a formal group endomorphism of $F(X,Y)$,
    i.e. $[n](F(X,Y)) = F ([n](X), [n](Y))$.
    \label{prop: nseries_end}
\end{proposition}

\begin{proof}
    We use induction on $n$ to prove it. For the base cases $n = 0,1$, the result follows immediately.
    Assume that the result holds for some natural number $k$. Now consider the case $k+1$.
    \begin{align*}
        [k+1]F(X,Y) &= F([k]F(X,Y), F(X,Y)) \\
            & = F(F([k](X), [k](Y)), F(Y,X)) \\
            & = F ([k](X), F([k](Y), F(Y,X))).
    \end{align*}
    Next, we focus on the term $F([k](Y), F(Y,X))$.
    \begin{equation*}
        F([k](Y), F(Y,X)) = F(F([k](Y), Y), X) = F([k+1](Y), X) = F(X,[k+1](Y)).
    \end{equation*}
    Hence, we have
    \begin{align*}
        [k+1]F(X,Y) & = F ([k](X), F(X,[k+1](Y))) \\
            & = F (F ([k](X), X), [k+1](Y)) \\
            & = F ([k+1](X), [k+1](Y)).
    \end{align*}
\end{proof}

\begin{remark}
    By induction, we deduce that $[n](X) = nX + O(X^2)$. As a result, for a ring $R$ of
    characteristic $p$, we have $[p](X) = a X ^ k + O(X ^ {k+1})$ for some $k> 1$.
\end{remark}

\begin{example}
    The $p$-series of additive formal group law is $[p]_{\GG_a} (X) = p X$.
    The $p$-series of multiplicative formal group law is $[p]_{\GG_m}(X) = (1+X)^p - 1$.
\end{example}

\begin{proposition}
    Assume that $R$ is of characteristic $p$ and nontrivial. Then the additive formal group law
    $\GG_a$ and the multiplicative formal group law $\GG_m$ over $R$ are not isomorphic.
\end{proposition}

\begin{proof}
    Assume that $\GG_a$ is isomorphic to $\GG_m$, then there is a formal group
    isomorphism $f(X)$ from $\GG_a$ to $\GG_m$.
    By the definition of homomorphism of formal group law, we have that

    \[ f([p]_{\GG_a}(X)) = [p]_{\GG_m}(f(X)). \]

    Note that

    \[[p]_{\GG_a}(X) = pX = 0 \quad \text{and} \quad [p]_{\GG_m} = (1+X) ^ p - 1 = X ^ p.\]

    Then we have

    \[f([p]_{\GG_a}(X)) = f(0) = 0 \quad \text{and} \quad [p]_{\GG_m}(f(X)) = (f(X))^p. \]

    Therefore, we have the equality
    \begin{equation}
        (f(X)) ^ p = 0.
        \label{eq: pow_zero}
    \end{equation}

    If we write
    $f(X) = \sum_{n=1}^{\infty} a_n X ^ n$, then $a_1$ is invertible in $R$, because $f$ is
    a formal group isomorphism. By equation \eqref{eq: pow_zero}, we have that $a_1 ^ p = 0$, which
    contradicts to $a_1$ is invertible in $R$.

\end{proof}

\begin{proposition}
    Assume $R$ is of characteristic $p$. Then $[p](X) = 0$ or
    $[p](X) = a X ^ {p^n} + O (X ^ {p^n + 1})$.
\end{proposition}

\begin{proof}
    The proof can be found in Corollary \ref{cor: p_series_form}.

    TODO: need more detailed.
\end{proof}

\begin{definition}[{\cite[Definition 1.5]{hazewinkel1978formal}}]
    Let $F(X,Y)= X + Y + \sum_{i, j \ge 1} c_{ij} X ^i Y^j$ be a formal group law over a ring $R$,
    and let $\phi : R \to R'$ be a ring homomorphism.
    By applying
    $\phi$ to the coefficients of $F(X,Y)$, we obtain a formal group law over $R'$
    \begin{equation}
        \phi_* F(X,Y) = X + Y + \sum_{i,j \ge 1} \phi (c_{ij}) X ^i Y^j.
    \end{equation}
    \label{defn: pushforward of coeff ring}
\end{definition}

\begin{remark}
    There is a functor from the category of commutative rings to \textbf{Set}, which maps a commutative ring
    $R$ to all formal group laws over $R$. Given a ring homomorphism $\phi : R \to S$, the definition
    above provides a map from formal group laws over $R$ to formal group laws over $S$.

\end{remark}
